\clearpage\pagestyle{otMain}\fontsize{11}{15}\selectfont
\corrigesChapitre{4}{Équations et inéquations}
\begin{blocLivre}{cor-4-1}
\begin{corrige}{1}{Une équation linéaire}
\demarche{Isoler l’inconnue par des opérations réversibles.}
On ajoute 7 aux deux membres : $3x=18$. On divise par 3, qui est non nul : $\res{x=6}$. La substitution donne $3\times6-7=11$. L’ensemble des solutions est $\res{S=\{6\}}$ ; il contient une seule valeur.
\end{corrige}
\end{blocLivre}
\begin{blocLivre}{cor-4-2}
\begin{corrige}{2}{L’inconnue des deux côtés}
\demarche{Développer, regrouper les termes en x puis les constantes.}
On développe : $5x-10=2x+8$. En retirant $2x$ puis en ajoutant 10, on obtient $3x=18$, donc $\res{S=\{6\}}$. Contrôle : à gauche $5(6-2)=20$ et à droite $12+8=20$. Les deux membres ont bien la même valeur.
\end{corrige}
\end{blocLivre}
\begin{blocLivre}{cor-4-3}
\begin{corrige}{3}{Un signe à renverser}
\demarche{Repérer la division par un nombre négatif.}
On soustrait 3 : $-4x>12$. En divisant par $-4$, on renverse le signe : $x<-3$. Ainsi $\res{S=]-\infty;-3[}$. La borne $-3$ est exclue puisque l’inégalité est stricte ; elle donnerait l’égalité 15.
\end{corrige}
\end{blocLivre}
\begin{blocLivre}{cor-4-4}
\begin{corrige}{4}{Un carré donné}
\demarche{Distinguer un second membre positif, nul ou négatif.}
Pour $x^2=49$, il faut les deux signes : $\res{S=\{-7;7\}}$. Pour $x^2=0$, la seule solution est $\res0$. Enfin un carré réel est positif ou nul, donc $x^2=-9$ n’a aucune solution réelle : $\res{S=\varnothing}$.
\end{corrige}
\end{blocLivre}
\begin{blocLivre}{cor-4-5}
\begin{corrige}{5}{Produit nul}
\demarche{Remplacer le produit nul par une alternative.}
Le produit est nul si $x-5=0$ ou $2x+3=0$. Ces deux équations donnent $x=5$ et $x=-3/2$. On prend leur réunion : $\res{S=\{-3/2;5\}}$. Il ne faut pas imposer que les deux facteurs s’annulent simultanément.
\end{corrige}
\end{blocLivre}
\begin{blocLivre}{cor-4-6}
\begin{corrige}{6}{Une équation quotient}
\demarche{Fixer le domaine puis supprimer le dénominateur.}
Le domaine impose $x\ne3$. Sur ce domaine, on peut multiplier par $x-3$ : $2x+1=5x-15$, donc $3x=16$. La solution $\res{x=16/3}$ est autorisée. Le dénominateur vaut $7/3$ et le numérateur $35/3$, dont le quotient est bien 5.
\end{corrige}
\end{blocLivre}
\begin{blocLivre}{cor-4-7}
\begin{corrige}{7}{Tout, rien ou une valeur}
\demarche{Lire l’égalité finale quand l’inconnue disparaît.}
La première devient $2x+2=2x+2$, toujours vraie : $\res{S=\mathbb R}$. La seconde devient $2=3$, impossible : $\res{S=\varnothing}$. L’annulation des termes en $x$ ne permet pas de conclure à $x=0$ : il faut examiner l’égalité restante.
\end{corrige}
\end{blocLivre}
\begin{blocLivre}{cor-4-8}
\begin{corrige}{8}{Un produit de deux facteurs}
\demarche{Découper la droite aux zéros de chaque facteur.}
Les zéros sont $-3$ et 1. Avant $-3$, les deux facteurs sont négatifs, donc le produit est positif. Entre $-3$ et 1, un seul est négatif : le produit est négatif. Après 1, il est positif. Les zéros sont inclus par $\le$ : $\res{S=[-3;1]}$.
\label{fig-cor-4-8}
\figureGeo[corrige]{signes-produit}
\end{corrige}
\end{blocLivre}
\begin{blocLivre}{cor-4-9}
\begin{corrige}{9}{Un quotient et sa valeur interdite}
\demarche{Séparer un zéro autorisé d’un zéro de dénominateur.}
Le quotient est défini pour $x\ne-1$ ; le numérateur s’annule en 2. Il est positif lorsque les deux facteurs ont le même signe : pour $x<-1$ et pour $x>2$. Les bornes sont exclues, $-1$ étant interdit et 2 donnant 0. Ainsi $\res{S=]-\infty;-1[\cup]2;+\infty[}$.
\label{fig-cor-4-9}
\figureGeo[corrige]{signes-quotient}
\end{corrige}
\end{blocLivre}
\begin{blocLivre}{cor-4-10}
\begin{corrige}{10}{Déplacer une expression}
\demarche{Comparer par différence plutôt que diviser par une expression inconnue.}
On ramène tout à gauche et on factorise le facteur commun : $(x-2)((x-2)-(x+4))\le0$, soit $-6(x-2)\le0$. La division par $-6$ renverse l’ordre : $x-2\ge0$. La réponse est $\res{S=[2;+\infty[}$. Diviser directement par $x-2$ aurait demandé de distinguer son signe.
\end{corrige}
\end{blocLivre}
\begin{blocLivre}{cor-4-11}
\begin{corrige}{11}{Un abonnement avantageux}
\demarche{Traduire « strictement moins cher » puis revenir aux entiers.}
Pour $n\in\mathbb N$, les coûts sont $8n$ et $30+5n$. L’abonnement est moins cher si $30+5n<8n$, donc $n>10$. Le premier entier qui convient est $\res{11}$. À 10 séances les prix sont égaux ; à 11, ils valent 88 euros et 85 euros.
\end{corrige}
\end{blocLivre}
\begin{blocLivre}{cor-4-12}
\begin{corrige}{12}{Un budget avec remise fixe}
\demarche{Écrire une inégalité de budget avec un domaine entier positif.}
Si $n$ est le nombre entier de carnets, $4{,}5n+7\le52$. Donc $4{,}5n\le45$, soit $n\le10$. On peut commander $\res{10\text{ carnets}}$. Le coût est alors exactement 52 euros ; onze carnets coûteraient 56,50 euros et dépasseraient le budget.
\end{corrige}
\end{blocLivre}
\begin{blocLivre}{cor-4-13}
\begin{corrige}{13}{Des dimensions à retrouver}
\demarche{Factoriser une expression simple puis filtrer par le contexte.}
On écrit $(x+3)(x-1)=21$, soit $x^2+2x-24=0$. On reconnaît $(x+6)(x-4)=0$, donc $x=-6$ ou $x=4$. La condition $x>1$ élimine $-6$. Ainsi $\res{x=4}$ ; les dimensions sont 7 cm et 3 cm, dont le produit vaut 21.
\label{fig-cor-4-13}
\figureGeo[corrige]{rectangle-equation}
\end{corrige}
\end{blocLivre}
\begin{blocLivre}{cor-4-14}
\begin{corrige}{14}{Un paramètre au premier degré}
\demarche{Repérer la valeur du paramètre qui annule le coefficient.}
Si $m\ne2$, le coefficient de $x$ est non nul et l’unique solution est $\res{x=6/(m-2)}$. Si $m=2$, l’équation devient $0x=6$, impossible : elle n’a aucune solution. La séparation de ces deux cas est nécessaire avant la division par $m-2$.
\end{corrige}
\end{blocLivre}
\begin{blocLivre}{cor-4-15}
\begin{corrige}{15}{Une équation qui perdrait zéro}
\demarche{Factoriser avant toute division par l’inconnue.}
On ramène tout à gauche : $x(x-3)-2x=x(x-5)=0$. Les solutions sont donc $\res{S=\{0;5\}}$. Diviser par $x$ suppose $x\ne0$ ; on ne trouverait alors que 5 et on aurait éliminé la solution 0. La vérification dans l’équation initiale confirme les deux valeurs.
\end{corrige}
\end{blocLivre}
